% ch09_t3b.tex: Chapter 9, Few-shot layout adaptation (track T3b, E1 opening #9).
% Every number is traced in ch09_t3b_NOTES.md. Sections 9.1 and 9.2 were written before the read; the result (read once on 23 Nov
% 2026) was filled in afterwards from verdict.txt and read.txt of the T3b read only (lines in the NOTES file, section 5).
\chapter{Few-shot layout adaptation: track T3b, E1 opening \#9}\label{ch:t3b}

Chapter~\ref{ch:e1_result} tested the final recipe zero-shot: no model had seen an E1 cell or a labelled half-cut cell. On the 62
half-cut cells of E1, a layout of which the training pool holds none, the gain of C2 over C1 was not detected (+0.0110
[$-$0.0289, +0.0722]). This chapter asks a second question of the same cells: how much does the same recipe gain on E1's layouts
once it has labelled cells of them? We call this the few-shot track (track T3b of the protocol, unrelated to
Table~\ref{tab:T3b}); it is E1 opening \#9. Sections~\ref{sec:t3b-design} and~\ref{sec:t3b-rule} give the design and the
registered rule, Section~\ref{sec:t3b-result} the result, and Section~\ref{sec:t3b-scope} what the track can and cannot show.

\section{Motivation and design}\label{sec:t3b-design}

\paragraph{Why.} The pool has no half-cut cell, so CV5 cannot measure that layout and the delivered model has never seen a labelled
cell of it. E1 holds 62 half-cut cells and 37 cells with a scuff streak of at least 500~px; the only other labelled half-cells of
this work, 36 E-SCDD cells, carry speck-level crack labels and serve for validation only. Our reading (judgement): what labelled
cells of a new layout buy is what WattCycle needs to know to plan labelling, and only E1 can measure it here.

\paragraph{Set and folds.} E1 holds 106 cells from 71 source images, one module each; 32 cells hold crack labels (32,294 crack
pixels, 296 crack components of at least 20~px) and 74 are crack-free. The hard stratum is the half-cut cells with a scuff streak;
its 20 crack cells hold 75\% of E1's crack pixels. We split E1 into $K=5$ folds by source image, drawn deterministically on 12
November from the E1 lists and ground-truth counts only (no prediction), with the crack pixels of the hard stratum balanced first,
and frozen with their md5 before any run. The held-out folds hold 22, 22, 20, 21 and 21 cells from 12, 14, 14, 16 and 15 images.

\paragraph{Arms.} For fold $k$, the arm ALL trains on the 532 pool cells plus every E1 cell outside fold $k$ (84 to 86 cells);
the arm HALF trains on the 532 plus a fixed half of those source images (42 or 43 cells), chosen deterministically to hold half of
the training side's cells and half of its hard-stratum crack pixels (shares 0.494 to 0.506 and 0.496 to 0.501). Each fold model
predicts only its held-out cells, and the five fold records of an arm and seed join into one out-of-fold record over the 106
cells. Three deployable models train on the 532 plus all 106 E1 cells and are never scored on E1. With seeds 5, 6 and 7, the
track has 33 trainings: 15 ALL, 15 HALF and 3 deployable.

\paragraph{Recipe, labels and inference.} Every model uses exactly the recipe of C2 (HRNet-W18 U-Net++
\citep{wang2021hrnet,zhou2018unetpp}, PF encoder schedule, SWAD \citep{cha2021swad} over steps 10,000 to 20,000, 20,000 steps, label
set v1.1 on the pool; Chapter~\ref{ch:methods}); only the training list differs. E1 labels enter in their scoring view
(Chapter~\ref{ch:protocol}), with classes 25 to 28 and ignore pixels set to ignore, so the training label equals the scoring label
on every scored pixel. The validation monitor is a fixed set of 107 pool training cells whose numbers are never read, and no fold's
training list lists or links its held-out cells. Prediction is C2's: the \texttt{swad} readout of the final weights, D4 with 8
views, argmax (OP-A). No threshold, bias, view weight or checkpoint was set on, or chosen with, any E1 cell. From staging to
clean-up, the training machine held the labels of all 106 cells, which the deployable models need; each fold's held-out labels were
present there but not linked from its training root.

\paragraph{The comparator.} ZS is the zero-shot C2 of opening \#8 for seeds 5, 6 and 7: the same recipe, seeds and inference,
never trained on E1. Its predictions were made at opening \#8 and are re-scored here. With 3 models, its level is not the five-seed
C2 level of Chapter~\ref{ch:e1_result}.

\section{Registration, rule and power}\label{sec:t3b-rule}

\paragraph{Registration and run.} The rule was written on 21 November, before any model of the track was trained; after an independent
review it was registered on 22 November at 03:13 IST, and the ledger row states that from that date E1 is training data for the few-shot models. All 33
trainings ended without error on 22 November, with no repair run. The opening ran on 23 November from 08:40:46 to 13:39:27 IST.
The 30 sealed prediction jobs, one per fold model, all ended without error. Of these, 29 ran each on a GPU that held no other job above
4,096~MiB. The GPUs were then busy with other users' jobs, so sharing was allowed for the last two jobs (13:34:35 to 13:35:49 IST),
outside the reviewed plan and logged at the time: one started on a GPU that had freed, and the last ran beside another user's job on
the same GPU. The dumps are deterministic fp32 inference on one GPU model (RTX A6000), and the read checks each dump's manifest. The
read ran once, after all 30 existed, from 13:38:00 to 13:39:27 IST, without error; it re-scored the three ZS dumps and found them
identical to their opening \#8 records.

\paragraph{Rule.} The metric is crack relF1@2 at OP-A, pooled over the scored pixels of the 106 cells, as at opening \#8. The
primary contrast P is $\Delta$, the mean over seeds 5 to 7 of the out-of-fold relF1@2 of ALL minus the same mean for ZS. Its
interval is the seed-aware two-stage bootstrap of opening \#8 (per draw the 71 clusters with replacement, then for each arm 3 models
with replacement from its 3; 2,000 draws; 95\% percentile), and $M=0.010$. P is CONFIRMED if $\Delta\geq M$, the lower end is
above 0, and two guards pass: G1, false crack pixels on the 74 crack-free cells, ALL at most $1.5\times$ ZS; G2, a change in
mIoU-R13 of at least $-$0.010. With a failing guard the verdict is GAIN, FAILING GUARDS NAMED; otherwise it is NOT CONFIRMED, worded
WORSE, ``interval holds 0'' or ``below M''. The width ratio ALL / ZS is printed but is not a guard: E1 labels are about $1.5\times$
wider than the pool labels, so a model trained on E1 cells learns the wider width. Everything else (HALF $-$ ZS, ALL $-$ HALF,
any-hit crack component recall, the strata, inactive IoU, strict crack IoU with its band $\pm$0.019) is printed beside and decides
nothing.

\paragraph{Power.} The rule stated its power before any number (judgement in the rule). Opening \#8's seed-aware standard error was
0.0265 with 5 models per arm; with 3, the rule expected about 0.030, giving power of about 52\% at a true +0.06, 76\% at +0.08 and
91\% at +0.10; the smallest effect with 80\% power is about +0.085. The rule judged +0.10 or more plausible if labels of the layout
close part of the gap on the hard stratum, where an earlier zero-shot model of this work recalled about 0.13 of the crack pixels.
The cell count overstates the information: half of E1's crack pixels lie in 3 source images (the top 1, 3, 5 and 10 hold 0.22,
0.50, 0.69 and 0.86), and the 20 hard crack cells come from 10 images. Two held-out folds are in effect one image each, so only the
pooled score is usable, with the uncertainty of about 9 images (effective $n$ 8.6). Smaller gains, ALL $-$ HALF and the strata are
under-powered, and inactive (10 cells from 7 images, 87\% of its pixels in 3) cannot be answered.

\section{Result}\label{sec:t3b-result}

The read printed: ``T3B VERDICT: P (ALL $-$ ZS pooled out-of-fold crack relF1@2) CONFIRMED; Delta +0.2213 [+0.0977, +0.3504];
width ratio ALL / ZS 1.426 (reported).'' It ran once on 23 November 2026, from 13:38:00 to 13:39:27 IST.

\paragraph{Primary contrast.} $\Delta$, the out-of-fold crack relF1@2 of ALL minus that of ZS at OP-A, each the mean over seeds 5
to 7, is +0.2213, with seed-aware 95\% interval [+0.0977, +0.3504] (SE 0.0643) and fixed-model interval [+0.1066, +0.3390]. It is
above $M=0.010$ and its lower end is above 0. Crack relF1@2 (mean over seeds 5 to 7, range over seeds in brackets) is 0.4881
[0.4656, 0.5067] for ZS, 0.6924 [0.6797, 0.7005] for HALF and 0.7094 [0.6943, 0.7309] for ALL; relP@2 is 0.6098, 0.8119 and
0.8403, and relR@2 0.4095, 0.6038 and 0.6140. Precision and recall both rise (ALL $-$ ZS: relP@2 +0.2306 [+0.1394, +0.3440],
relR@2 +0.2046 [+0.0356, +0.3835]).

\paragraph{Guards and width.} Both guards pass. G1: on the 74 crack-free cells ALL predicts 541 false crack pixels against 2231
for ZS (ratio 0.242; bound 1.5). G2: mIoU-R13 changes by +0.1385 (bound $-$0.010). The width ratio ALL / ZS is 1.426 (ALL 0.631,
ZS 0.443), printed and not a guard. E1 labels are about $1.5\times$ wider than the pool labels, and a model trained on E1 cells
learns that width, which lifts strict IoU and pixel recall without finding more cracks.

\paragraph{Printed beside.} HALF $-$ ZS is +0.2042 [+0.0718, +0.3366] (fixed-model [+0.0773, +0.3274]), about the same as ALL $-$
ZS, and ALL $-$ HALF is +0.0171 [$-$0.0246, +0.0625] (fixed-model [$-$0.0183, +0.0506]), not separated: half of the labelled
training-side images gave about the same gain as all of them. Of the 296 crack components of at least 20~px, any-hit recall (a
component counts if any of its pixels lies within 2~px of a predicted crack pixel; arm mean over seeds, [min, max]) is 0.400
[0.301, 0.490] for ZS, 0.682 [0.645, 0.703] for HALF and 0.666 [0.635, 0.693] for ALL; recall at $\geq$50\% coverage is 0.245,
0.591 and 0.598, and component precision 0.659, 0.817 and 0.836. Our reading (judgement): a rise in component recall of this size is
hard to explain by line width alone, so the gain is not only the learned width.

\paragraph{Strata.} Table~\ref{tab:t3b-strata} gives the strata, which decide nothing. The half-cut cells carry the gain
(+0.2530 [+0.0878, +0.3885]), most of all the hard stratum, half-cut cells with scuff (+0.3002 [+0.1445, +0.4748]). On half-cut
cells without scuff the gain is not separated and its point estimate is negative ($-$0.0701 [$-$0.3019, +0.0517]); there ZS already
scores 0.8870 [0.8797, 0.8928], against 0.8168 [0.8034, 0.8340] for ALL. The full cells are separated but wide (+0.1525 [+0.0260,
+0.4423]); on them HALF $-$ ZS is not separated (+0.0818 [$-$0.0433, +0.2218]), and ALL $-$ HALF is +0.0708 [$-$0.0019, +0.3419],
against +0.0064 [$-$0.0484, +0.0629] on the half-cut cells. On the 74 crack-free cells every arm scores relF1@2 0.0000, with 541
(ALL), 1366 (HALF) and 2231 (ZS) false crack pixels. The CSIR cells repeat the pooled result. The BMRK stratum (5 cells, 5 clusters)
is degenerate: ALL and ZS both score 0.0000 there and HALF 0.5443 with a seed range of [0.0000, 0.9908], so it says nothing.

\begin{table}[htbp]
\centering
\scriptsize
\setlength{\tabcolsep}{3pt}
\caption{Few-shot track (E1 opening \#9): crack relF1@2 contrasts at OP-A, out of fold, mean over seeds 5 to 7, with seed-aware
95\% intervals, for all cells and per stratum. Only ALL $-$ ZS on all cells decides; the strata decide nothing. Hard = half-cut
cells with scuff. The 74 crack-free cells score 0 for every arm and are left out.}
\label{tab:t3b-strata}
\begin{tabular}{@{}>{\raggedright\arraybackslash}p{1.8cm}lccc@{}}
\toprule
Cells & Cells / clusters & ALL $-$ ZS & HALF $-$ ZS & ALL $-$ HALF \\
\midrule
All (P) & 106 / 71 & +0.2213 [+0.0977, +0.3504] & +0.2042 [+0.0718, +0.3366] & +0.0171 [$-$0.0246, +0.0625] \\
Full & 44 / 24 & +0.1525 [+0.0260, +0.4423] & +0.0818 [$-$0.0433, +0.2218] & +0.0708 [$-$0.0019, +0.3419] \\
Half-cut & 62 / 47 & +0.2530 [+0.0878, +0.3885] & +0.2466 [+0.0727, +0.3976] & +0.0064 [$-$0.0484, +0.0629] \\
Hard & 29 / 17 & +0.3002 [+0.1445, +0.4748] & +0.2929 [+0.1201, +0.4835] & +0.0073 [$-$0.0609, +0.0641] \\
Half-cut without scuff & 33 / 31 & $-$0.0701 [$-$0.3019, +0.0517] & $-$0.0704 [$-$0.2899, +0.0424] & +0.0003 [$-$0.0506, +0.0516] \\
Crack cells & 32 / 21 & +0.1967 [+0.0680, +0.3222] & +0.1927 [+0.0597, +0.3243] & +0.0040 [$-$0.0421, +0.0473] \\
CSIR & 101 / 66 & +0.2215 [+0.0977, +0.3415] & +0.2038 [+0.0708, +0.3354] & +0.0177 [$-$0.0267, +0.0647] \\
BMRK (degenerate) & 5 / 5 & +0.0000 [+0.0000, +0.0000] & +0.5443 [+0.0000, +0.8745] & $-$0.5443 [$-$0.8745, +0.0000] \\
\bottomrule
\end{tabular}
\end{table}

\paragraph{Other metrics.} Inactive IoU ALL $-$ ZS is +0.1463 [$-$0.0525, +0.4829], reported only (10 inactive cells in 7
images). Strict crack IoU ALL $-$ ZS is +0.1803 [+0.0929, +0.2606], outside the convention band $\pm$0.019; it decides nothing,
and the learned label width lifts it. Per class (arm means; decide nothing), crack IoU is 0.1866 for ZS, 0.3438 for HALF and
0.3669 for ALL; the largest change is sp mono, 0.1924, 0.2039 and 0.7966, whose cause was not examined; material and edge dark
score 0.0000 for every arm, and splice is not scored. Figure~\ref{fig:F20} shows the contrasts and levels per stratum.

\FIG{F20}{Few-shot track (E1 opening \#9). (a) Crack relF1@2 contrasts ALL $-$ ZS, HALF $-$ ZS and ALL $-$ HALF, with seed-aware
95\% intervals, for all cells and per stratum; numbers in brackets are cells / clusters; hard = half-cut cells with scuff; solid
line: zero; dashed line: M = 0.010. (b) Crack relF1@2 levels of ALL, HALF and ZS, as the mean over the three seeds; the bar runs from
the lowest to the highest seed and is not an interval. ALL is trained on the pool plus the E1 cells of the other 4 folds; HALF on
the pool plus a fixed half of those image ids. Only P = ALL $-$ ZS on all cells decides (confirmed, +0.2213 [+0.0977, +0.3504]);
HALF $-$ ZS and ALL $-$ HALF are printed beside and decide nothing. On half-cut cells without scuff, the ZS level (0.8870) is above
those of ALL and HALF. BMRK (5 cells; ALL and ZS score 0) and the 74 crack-free cells (every arm scores 0) are degenerate and not
drawn. E1, out of fold over 5 folds of E1, seeds 5 to 7, OP-A; ZS = zero-shot C2 seeds 5 to 7; 106 cells, 71 clusters; strata
decide nothing.}

\section{What the track can and cannot show}\label{sec:t3b-scope}

\paragraph{Can.} The track can show whether labelled cells of E1's layouts, scored out of fold by source image, raise the crack
relF1@2 of the C2 recipe on E1 by about +0.085 or more against the same recipe zero-shot, without more false crack pixels than the
guard allows. Our reading of the result (judgement): labelled cells of E1's layouts raise the crack relF1@2 of the C2 recipe
strongly against the same recipe zero-shot (+0.2213 [+0.0977, +0.3504], out of fold inside E1), and half of the labelled images
(42 or 43 E1 cells per fold model) gave about the same gain as all of them. For planning labels at WattCycle, on this evidence a
first batch of labelled cells of a new layout is worth having, and a second batch of the same size added nothing this track can
detect. It does not say how few cells suffice, and it does not show that the model handles a new layout without labels.

\paragraph{Cannot.} The rule had 80\% power only for gains of about +0.085 or more; the interval of $\Delta$, [+0.0977, +0.3504],
is wide, so the size of the gain is not pinned down. A training side holds 14 to 18 hard crack cells from at most 9 images, so the track has two points of a learning curve
(about half, and all), and which images fall in the half matters more than how many. Folds by source image may be optimistic for a
new site: E1's CSIR image ids run from 300 to 2,328, and whether neighbouring ids are modules of one batch with the same handling
marks was not measured. relF1@2 is less sensitive than strict IoU to the wider E1 labels but not immune, which is why any-hit
component recall is printed. 101 of the 106 cells are CSIR, and each arm has 3 seeds. The deployable models have no held-out score;
in our judgement their expected score on such cells is the out-of-fold score of ALL.

\paragraph{Cost.} From the registration of opening \#9, E1 is training data for the few-shot models. Every such run and record
carries a \texttt{trained\_on\_e1} mark that the read tools refuse outside the few-shot read; no number of this chapter enters the
zero-shot tables of Chapter~\ref{ch:e1_result} or an E1 column of any other track, and none supports a claim about unseen modules
without labels. The zero-shot result of Chapter~\ref{ch:e1_result} is therefore the last zero-shot result on E1 in this work, and a
later unseen-module claim needs a new labelled set (Chapter~\ref{ch:limitations}).
